
\section{Métodos autônomos de planejamento de trajetória}

Os métodos autônomos de guiagem (\emph{Autonomous Guidance Methods  (AGMs)}), têm sido estudados desde 2016 para melhorar a autonomia de voo em diferentes situações, esses métodos geram comandos de guiagem em tempo real para que a trajetória atenda às restrições dinâmicas do veículo durante todas as fases do voo.
Essa abordagem é diferente das tradicionais que são trajetórias pré-planejadas, o diferencial é que os AGMs lidam melhor com restrições não lineares e variáveis durante o voo \cite{song2023}.

Inicialmente, dentre os métodos existentes, destacam-se abordagens como a \emph{Computational Guidance}, onde as leis de controle são substituídas por algoritmos baseados em modelos de redes neurais.
Esse tipo de método elimina a necessidade de ajustes prévios e planejamento extensivo.
Outro modelo relevante é o \emph{Model-Based Real-Time Optimization}, que otimiza a trajetória de forma contínua, lidando com restrições complexas \cite{song2023}. Por fim, a \emph{Autonomous Mission Planning} detém o maior nível de autonomia e permite que o próprio sistema gerencie a missão em tempo real.

A aplicação desses métodos para missões de RVD/B exige considerações diferenciadas quando comparados a sistemas autônomos de navegação terrestre ou aeronáutica, pois no caso de veículos de lançamento, a não-linearidade da dinâmica de voo e a variação significativa da massa ao longo da missão impõem desafios adicionais para a alimentação e tomadas de decisões autônomas, por sua vez, os AGMs têm demonstrado capacidade de solucionar esses desafios ao longo do voo \cite{song2023}.